\uuid{myBw}
\exo7id{1011}
\titre{exo7 1011}
\auteur{ridde}
\organisation{exo7}
\datecreate{1999-11-01}
\isIndication{false}
\isCorrection{false}
\chapitre{Espace vectoriel}
\sousChapitre{Base}
\cpge{mpsi-mp2i.espaces-vectoriels.d, pcsi.espaces-vectoriels.d, ptsi.espaces-vectoriels.d, mpsi-mp2i.dimension-finie.b, pcsi.dimension-finie.b, ptsi.dimension-finie.b, bcpst1.espace-vectoriel-kn.a, bcpst1.espace-vectoriel-kn.b}
\module{Algèbre}
\niveau{L1}
\difficulte{}

\contenu{
\question{
On consid\`ere dans $\Rr^3$,  $\Pi = \text{vect}\left\{ (1, 1, 1), (1, 1, -1)\right\}$
et $D = \text{vect} \left\{ (0, 1, -1)\right\}$. Montrer que $\Rr^3 = \Pi \oplus D$.
}
}
